\documentclass[aps,prl,reprint,twocolumn,superscriptaddress,nofootinbib,floatfix]{revtex4-2}

\usepackage[utf8]{inputenc}
\usepackage[T1]{fontenc}
\usepackage{amsmath,amssymb,amsfonts,mathtools}
\usepackage{graphicx}
\usepackage{microtype}
\usepackage{booktabs}
\usepackage{enumitem}
\usepackage[dvipsnames]{xcolor}
\usepackage{silence}
\WarningFilter{nameref}{The definition of \label has changed!}
\usepackage{hyperref}

\hypersetup{
  colorlinks=true,
  linkcolor=NavyBlue,
  citecolor=Mahogany,
  urlcolor=TealBlue
}

\providecommand{\R}{\mathbb{R}}
\newcommand{\C}{\mathbb{C}}
\providecommand{\Sph}{\mathbb{S}}
\newcommand{\Tr}{\operatorname{Tr}}
\newcommand{\dif}{\,\mathrm{d}}
\providecommand{\Area}{\operatorname{Area}}
\newcommand{\norm}[1]{\left\|#1\right\|}
\DeclareMathOperator*{\argmin}{argmin}
\providecommand{\Hyp}{\mathbb{H}}
\providecommand{\II}{\mathrm{II}}
\providecommand{\Hsn}{\mathcal{H}}
\providecommand{\BV}{\mathrm{BV}}
\providecommand{\TV}{\mathrm{TV}}
\providecommand{\cutnorm}[1]{\|#1\|_{\square}}
\providecommand{\doi}[1]{\href{https://doi.org/#1}{\nolinkurl{doi:#1}}}

\begin{document}

\title{Observational Signatures of a Running Spectral Dimension:\texorpdfstring{\\}{ }Parametrization and Order-of-Magnitude Assessment}

\author{Reinaldo Maia Silva-Filho}
\email{reinaldo.filho1@estudante.ufla.br}
\affiliation{Programa de P\'os-Gradua\c{c}\~ao em Estat\'istica e Experimenta\c{c}\~ao Agropecu\'aria (PPGEE/DES), Universidade Federal de Lavras (UFLA), Lavras, MG, Brazil}

\date{\today}

\begin{abstract}
A spectral dimension that runs from $d_s = 4$ at large scales to $d_s = 2$ at short scales appears in several approaches to quantum gravity. For diffusion with the two-scale symbol $k^2 + \ell_P^2 k^4$ in four dimensions we give the return probability and $d_s(\tau)$ in closed form, and verify both against numerical quadrature. We then examine four phenomenological ans\"atze discussed in the companion chapters, none of which follows from the diffusion model: (i) a quartic correction to the graviton dispersion relation, $\omega^2 = c^2k^2(1 + \xi\ell_\ast^2 k^2)$; (ii) a running of the primordial tensor tilt postulated to track the spectral dimension, $\alpha_t = \frac{1}{2}(d_s - 4)$; (iii) analog measurements in quantum simulators; and (iv) atom interferometry. With $\ell_\ast = \ell_P$, the arrival-time delay between $10$\,Hz and $1$\,kHz for sources at $z = 1$--$8$ is $\sim 10^{-61}$--$10^{-60}$\,s. Evaluated at the physical momentum of the tensor modes at horizon exit, the tilt running is $|\alpha_t| \approx (H_{\mathrm{inf}}/M_P)^2 \lesssim 10^{-11}$. Both are many orders of magnitude below experimental reach. A detectable dispersion delay would require $\ell_\ast \sim 3\times 10^{-7}$\,m, a scale not motivated by the model. Analog simulators and atom interferometers do not test the framework without further input. The purpose of this letter is to record these estimates explicitly, so that the observational status of the framework is clear.
\end{abstract}

\maketitle

\emph{Introduction.}---The Planck scale $M_P = \sqrt{\hbar c / G} \approx 1.22 \times 10^{19}$\,GeV lies far beyond direct experimental reach, and perturbative quantum general relativity is non-renormalizable \cite{goroff1986ultraviolet}. Proposed signatures of quantum gravity must therefore rely on effects that accumulate over large distances or that are enhanced by some mechanism.

The companion chapters \cite{silvafilho2026grand} discuss a framework combining continuous tensor networks \cite{silvafilho2026functional, silvafilho2026flows, silvafilho2026spacetime}, fractional calculus on continuous Pascal simplices \cite{silvafilho2026pascal}, and $L^\infty$-minimax extrinsic curvature \cite{silvafilho2026noneuclidean}. Chapter~12 records a closed form for the spectral dimension of a diffusion with an assumed two-scale symbol, an elementary integral. Here we ask which observational signatures could be associated with it, and how large they are.

\begin{figure*}[t]
\centering
\includegraphics[width=\textwidth]{fig_experimental_signatures.pdf}
\caption{\textbf{Order-of-magnitude assessment.} (a) Arrival-time delay $|\Delta t_{\mathrm{disp}}|$ relative to $10$\,Hz, Eq.~\eqref{eq:delta_t} with $\xi = 1/2$. Blue: $\ell_\ast = \ell_P$ at $z = 1, 3, 8$. Red: the value $\ell_\ast \approx 3.0 \times 10^{-7}$\,m that would be required to reach a $0.1$\,ms delay at $1$\,kHz for $z = 3$. Dash-dotted: $0.1$\,ms timing benchmark. (b) The tilt-running ansatz \eqref{eq:tensor_running} as a function of $p/M_\ast$, with $p$ the physical momentum of the mode at horizon exit. The shaded region, $p = H_{\mathrm{inf}} \le 4.7 \times 10^{13}$\,GeV, contains the CMB modes for $M_\ast = M_P$ and $r < 0.036$. The script that generates the figure is in the accompanying code repository.}
\label{fig:experimental}
\end{figure*}

\emph{Running spectral dimension.}---The diffusion kernel with the isotropic symbol $k^2 + \ell_P^2 k^4$ in four Euclidean dimensions interpolates between the Laplacian at large scales and the bi-Laplacian at short scales, so $d_s$ runs from $4$ to $4/2 = 2$. The same end values occur in Causal Dynamical Triangulations \cite{ambjorn2005spectral}. In Ho\v{r}ava--Lifshitz gravity they arise from a different, anisotropic operator: with scaling exponent $z$ in $3+1$ dimensions the ultraviolet value is $d_s = 1 + 3/z$, so $d_s = 2$ needs $z = 3$, while $z = 2$ (a spatial $k^4$ term) gives $5/2$ \cite{horava2009spectral}. In $2+1$ dimensions, Sotiriou, Visser and Weinfurtner showed that Ho\v{r}ava--Lifshitz-type dispersion relations reproduce the CDT flow of $d_s$ \cite{sotiriou2011spectral}. Writing $u = k^2$ and $z \coloneqq \sqrt{\tau}/(2\ell_P)$, the return probability is
\begin{align}
    \label{eq:heat_kernel_exact}
    P(\tau) &= \frac{1}{16\pi^2} \int_0^\infty u\, e^{-\tau u - \tau \ell_P^2 u^2} \dif u \nonumber \\
    &= \frac{1}{32\pi^2 \ell_P^2 \tau} \Big[ 1 - \sqrt{\pi}\, z\, e^{z^2} \operatorname{erfc}(z) \Big],
\end{align}
and $d_s(\tau) \coloneqq -2\, \dif \ln P / \dif \ln \tau$ evaluates to
\begin{equation}
    \label{eq:ds_tau}
    d_s(\tau) = 1 - \frac{\tau}{2\ell_P^2} + \left[ 1 - \frac{\sqrt{\pi \tau}}{2\ell_P} e^{\tau/4\ell_P^2} \operatorname{erfc}\!\left( \frac{\sqrt{\tau}}{2\ell_P} \right) \right]^{-1}.
\end{equation}
Using $\operatorname{erfc}(z) = 1 - 2z/\sqrt{\pi} + \mathcal{O}(z^3)$ as $z \to 0$ and $\sqrt{\pi}\, z\, e^{z^2} \operatorname{erfc}(z) = 1 - 1/(2z^2) + 3/(4z^4) + \mathcal{O}(z^{-6})$ as $z \to \infty$,
\begin{equation}
    \label{eq:ds_limits}
    \lim_{\tau \to 0} d_s(\tau) = 2, \qquad \lim_{\tau \to \infty} d_s(\tau) = 4.
\end{equation}
In the infrared limit the explicit $-2z^2$ term cancels against the $2z^2 + 3 + \mathcal{O}(z^{-2})$ tail of the reciprocal bracket, and $d_s = 4 - 12\ell_P^2/\tau + \mathcal{O}(\ell_P^4/\tau^2)$. Eqs.~\eqref{eq:heat_kernel_exact} and \eqref{eq:ds_tau} agree with direct numerical quadrature to better than $10^{-6}$ for $\tau/\ell_P^2 \in [10^{-3}, 10^2]$.

\emph{Graviton dispersion.}---The running of $d_s$ does not by itself fix a dispersion relation for gravitational waves. As a phenomenological ansatz, we consider
\begin{equation}
    \label{eq:dispersion}
    \omega^2(k) = c^2 k^2 \left( 1 + \xi \ell_\ast^2 k^2 \right),
\end{equation}
where the diffusion model suggests the scale $\ell_\ast = \ell_P$ but does not fix $\xi$. Continued to Lorentzian signature, the isotropic symbol $p^2(1 + \ell_P^2 p^2)$ vanishes on $p^2 = 0$, a massless branch without dispersion ($\xi = 0$), and on $p^2 = -1/\ell_P^2$, a massive ghost. The anisotropic operator $\omega^2 + c^2(k^2 + \ell_P^2 k^4)$, which would give $\xi = 1$, has $d_s \to 5/2$ rather than $2$ in the ultraviolet. We treat $\xi$ as a free parameter of order unity; all delays are linear in $\xi$, and we display $\xi = 1/2$. The group velocity is $v_g \approx c(1 + \frac{3}{2}\xi \ell_\ast^2 k^2)$. A graviton observed at frequency $f$ had physical wavenumber $2\pi f (1+z')/c$ at redshift $z'$, so the velocity excess is weighted by $(1+z')^2$ along the path \cite{jacob2008lorentz, mirshekari2012constraining}. Two frequencies emitted together at redshift $z$ therefore arrive separated by
\begin{equation}
    \label{eq:delta_t}
    |\Delta t_{\mathrm{disp}}| = \frac{6 \pi^2 \xi \ell_\ast^2}{c^3} D_2(z) \left( f_2^2 - f_1^2 \right),
\end{equation}
with $D_2(z) = \frac{c}{H_0} \int_0^z \frac{(1+z')^2\,\dif z'}{\sqrt{\Omega_m(1+z')^3 + \Omega_\Lambda}}$.
For $\ell_\ast = \ell_P$, $\xi = 1/2$, $f_1 = 10$\,Hz, $f_2 = 1$\,kHz and $H_0 = 70$\,km\,s$^{-1}$\,Mpc$^{-1}$, $\Omega_m = 0.3$, one finds $|\Delta t_{\mathrm{disp}}| = 6 \times 10^{-62}$, $3 \times 10^{-61}$ and $1.2 \times 10^{-60}$\,s at $z = 1, 3, 8$ [Fig.~\ref{fig:experimental}(a)]. As a benchmark for third-generation detectors such as the Einstein Telescope \cite{punturo2010einstein} and Cosmic Explorer \cite{reitze2019cosmic} we take a timing accuracy of $10^{-4}$\,s, an optimistic figure corresponding to a phase accuracy of $\sim 0.1$\,rad at $\sim 100$\,Hz for loud events; actual dispersion tests fit the waveform phase rather than arrival times. The Planck-scale signal is therefore $56$--$57$ orders of magnitude below this benchmark, and LISA \cite{amaro2017laser}, which observes millihertz frequencies, is less sensitive to an $f^2$ effect. A $10^{-4}$\,s delay at $z = 3$ would require $\ell_\ast \approx 3.0 \times 10^{-7}$\,m, corresponding to an energy scale $\hbar c/\ell_\ast \approx 0.7$\,eV. Nothing in the framework points to such a scale, and it would have to be confronted with existing constraints on gravitational-wave propagation, among them the bound $|v_g - c|/c \lesssim 10^{-15}$ from GW170817 \cite{abbott2017gw170817}.

\emph{Tensor-tilt running.}---In single-field slow-roll inflation, $\mathcal{P}_t(k) = A_t (k/k_*)^{n_t}$ with $n_t = -r/8$. A two-point Pad\'e form with the limits \eqref{eq:ds_limits} and an infrared approach of the same power as Eq.~\eqref{eq:ds_tau} (the coefficient depends on the identification $\tau \sim p^{-2}$ and is not matched) is $d_s(p) = 2 + 2/(1 + (p/M_\ast)^2)$. If one \emph{postulates} that the running of the tilt tracks the spectral dimension,
\begin{equation}
    \label{eq:tensor_running}
    \alpha_t \coloneqq \frac{\dif n_t}{\dif \ln k} = \frac{1}{2} \left( d_s(p) - 4 \right) = -\frac{(p/M_\ast)^2}{1 + (p/M_\ast)^2},
\end{equation}
the momentum $p$ must be a physical one. A local effect with scale $M_\ast$ acts on a tensor mode when the mode is generated, at horizon exit, where its physical momentum is $p = k/a = H_{\mathrm{inf}}$; the comoving wavenumber measured today depends on the normalization of the scale factor and is not an invariant scale. The bound $r < 0.036$ \cite{bicepkeck2021constraints} and $A_s \approx 2.1 \times 10^{-9}$ \cite{planck2018parameters} give, through $2H_{\mathrm{inf}}^2/(\pi^2 M_{\mathrm{red}}^2) = rA_s$ with $M_{\mathrm{red}} = M_P/\sqrt{8\pi}$, the bound $H_{\mathrm{inf}} \le 4.7 \times 10^{13}$\,GeV. For $M_\ast = M_P$ the running is then $|\alpha_t| \approx (H_{\mathrm{inf}}/M_P)^2 \le 1.5 \times 10^{-11}$ ($3.7 \times 10^{-10}$ for $M_\ast = M_{\mathrm{red}}$), nearly the same for all CMB multipoles because $H_{\mathrm{inf}}$ changes little while they exit the horizon [Fig.~\ref{fig:experimental}(b)]. LiteBIRD \cite{hazumi2019litebird} and CMB-S4 \cite{abazajian2016cmb} target the tensor amplitude at the level $r \sim 10^{-3}$; a running of $10^{-11}$ in the tensor tilt is far beyond their reach, and the ansatz predicts no observable modification of $C_\ell^{BB}$.

\emph{Quantum simulators.}---Programmable Rydberg arrays \cite{ebadi2021quantum} and superconducting processors can prepare tensor-network states, estimate their quantum Fisher metrics, and measure out-of-time-order correlators $F(t) = \langle W(t) V(0) W(t) V(0) \rangle_\beta$. The Lyapunov exponent obeys the Maldacena--Shenker--Stanford bound $\lambda_L \le 2\pi k_B T/\hbar$ \cite{maldacena2016bound}, which the SYK model approaches only at strong coupling. Such experiments probe the specific many-body systems realized. They would test the framework only if a concrete model were identified in which its conjectured relations (Conjecture~7.1 of \cite{silvafilho2026grand}) make a quantitative prediction. We do not have such a model.

\begin{table*}[t]
\centering
\caption{\textbf{Summary of the order-of-magnitude assessment} ($\ell_\ast = \ell_P$, $M_\ast = M_P$ unless stated).}
\label{tab:experiments}
\begin{ruledtabular}
\begin{tabular}{lllll}
\textbf{Frontier} & \textbf{Observable} & \textbf{Model value} & \textbf{Instrument reach} & \textbf{Status} \\
\midrule
Gravitational waves & \parbox[t]{3.0cm}{\raggedright $|\Delta t_{\mathrm{disp}}|$, $10$\,Hz--$1$\,kHz, $z = 1$--$8$} & $10^{-61}$--$10^{-60}$\,s & \parbox[t]{3.0cm}{\raggedright ET, CE: $10^{-4}$\,s (benchmark)} & \parbox[t]{3.0cm}{\raggedright $56$--$57$ orders below reach} \\
\addlinespace[3pt]
CMB & \parbox[t]{3.0cm}{\raggedright tilt running $\alpha_t$ at horizon exit} & $|\alpha_t| \lesssim 10^{-11}$ & \parbox[t]{3.0cm}{\raggedright LiteBIRD, CMB-S4: $r \sim 10^{-3}$} & \parbox[t]{3.0cm}{\raggedright unobservable} \\
\addlinespace[3pt]
Quantum simulators & \parbox[t]{3.0cm}{\raggedright Fisher metric, OTOC} & \parbox[t]{2.6cm}{\raggedright no model-specific prediction} & \parbox[t]{3.0cm}{\raggedright Rydberg arrays, transmons} & \parbox[t]{3.0cm}{\raggedright requires a concrete model} \\
\addlinespace[3pt]
Atom interferometry & \parbox[t]{3.0cm}{\raggedright $\ell_P/\lambda_{\mathrm{dB}}$, Sr, $v \lesssim 45$\,m/s} & $\lesssim 2 \times 10^{-25}$ & \parbox[t]{3.0cm}{\raggedright MAGIS-100, AION} & \parbox[t]{3.0cm}{\raggedright no test proposed} \\
\end{tabular}
\end{ruledtabular}
\end{table*}

\emph{Atom interferometry.}---Long-baseline atom interferometers (MAGIS-100 \cite{abe2021matter}, AION \cite{badurina2020aion}) do not test a ``Jordan-loop protection'' of loop-quantum-gravity holonomies. Such a test would require the restriction to non-self-intersecting loops to follow from chronology protection. Chronology constrains timelike curves, whereas the loops in question are spacelike, and the restriction is a conjecture about a kinematical sector (Conjecture~4.3 of \cite{silvafilho2026grand}). Moreover, for $^{87}$Sr atoms $\ell_P/\lambda_{\mathrm{dB}} = m v \ell_P/(2\pi\hbar) \approx 3.5 \times 10^{-27}\,(v/1\,\mathrm{m\,s^{-1}})$, which stays below $2 \times 10^{-25}$ for the launch velocities ($v \lesssim 45$\,m\,s$^{-1}$) of long-baseline instruments; no prediction of the framework is attached to this ratio.

\emph{Conclusion.}---The closed form \eqref{eq:ds_tau} gives an explicit, numerically verified interpolation of the spectral dimension between $2$ and $4$ for a two-scale diffusion. The observational ans\"atze associated with it, with the scale set by $\ell_P$, give effects suppressed by factors of order $(\ell_P k)^2$ for gravitational waves and $(H_{\mathrm{inf}}/M_P)^2$ for the tensor tilt, far below the reach of gravitational-wave and CMB experiments. For quantum simulators and atom interferometers the framework makes no prediction to compare with. The framework in its present form makes no observable prediction. Obtaining one would require an independent mechanism that enhances Planck-suppressed effects, or a lower scale $\ell_\ast$ justified from first principles.

\begin{acknowledgments}
This research was financed in part by the Coordena\c{c}\~ao de Aperfei\c{c}oamento de Pessoal de N\'ivel Superior -- Brasil (CAPES) -- Finance Code 001. The author thanks the faculty and researchers at UFLA PPGEE/DES for institutional support and valuable discussions.
\end{acknowledgments}

% Shared declarations block, \input by every chapter before its bibliography.
% Keep in sync with the "Declarations" section of master_book_unified_quantum_gravity.tex.
\section*{Declarations}

\noindent\textbf{Funding.} This research was financed in part by the Coordena\c{c}\~ao de Aperfei\c{c}oamento de Pessoal de N\'ivel Superior -- Brasil (CAPES) -- Finance Code 001.

\medskip\noindent\textbf{Competing interests.} The author declares no competing interests.

\medskip\noindent\textbf{Use of generative AI tools.} Large language model assistants (Anthropic Claude; Google Gemini, used through the Antigravity agent environment) were used extensively in preparing this work: drafting and editing text and \LaTeX{} source; proposing, expanding and revising derivations and proofs under the author's direction; writing the Python verification scripts and the \textsc{Lean 4} files; searching and checking the literature and references; and adversarial auditing of proofs, numerical claims and code. These audits found errors, some of them originally introduced with AI assistance. The errors and their corrections are recorded in the repository file \texttt{WORKPLAN.md}. AI tools are not authors. The author chose the research questions and the overall framework, reviewed the material, and is responsible for the content, including any remaining errors.

\medskip\noindent\textbf{Verification status.} Results labelled Theorem, Proposition or Lemma are proved in the text or cited as classical; statements labelled Conjecture are not proved. The numerical scripts compare formulas of the text with independent computations and are checks, not proofs. The \textsc{Lean 4} modules in \texttt{formal\_proofs\_book/Book} are a naming skeleton and do not verify the mathematics of this chapter.

\medskip\noindent\textbf{Code and data availability.} Sources, numerical scripts and the audit log are available at
\begin{center}\url{https://github.com/reinaldomsilvafilho-netizen/quantum-gravity-lean4}.\end{center}
The monograph is archived on Zenodo, \href{https://doi.org/10.5281/zenodo.22290043}{doi:10.5281/zenodo.22290043}, under the CC~BY~4.0 license.


\begin{thebibliography}{99}

\bibitem{abbott2017gw170817}
B.~P. Abbott et~al. (LIGO Scientific Collaboration, Virgo Collaboration, Fermi GBM, INTEGRAL),
\emph{Gravitational waves and gamma-rays from a binary neutron star merger: GW170817 and GRB 170817A},
Astrophys. J. Lett., \textbf{848} (2017), L13,
\doi{10.3847/2041-8213/aa920c}.

\bibitem{horava2009spectral}
P.~Ho\v{r}ava,
\emph{Spectral dimension of the universe in quantum gravity at a Lifshitz point},
Phys. Rev. Lett., \textbf{102} (2009), 161301,
\doi{10.1103/PhysRevLett.102.161301}.

\bibitem{sotiriou2011spectral}
T.~P. Sotiriou, M.~Visser, and S.~Weinfurtner,
\emph{Spectral dimension as a probe of the ultraviolet continuum regime of causal dynamical triangulations},
Phys. Rev. Lett., \textbf{107} (2011), 131303,
\doi{10.1103/PhysRevLett.107.131303}.

\bibitem{abazajian2016cmb}
K.~N. Abazajian et~al.,
\emph{CMB-S4 Science Book, First Edition},
arXiv:1610.02743 (2016),
\doi{10.48550/arXiv.1610.02743}.

\bibitem{abe2021matter}
M.~Abe et~al.,
\emph{Matter-wave atomic gradiometer interferometric sensor (MAGIS-100)},
Quantum Sci. Technol., \textbf{6} (2021), 044003,
\doi{10.1088/2058-9565/abf719}.

\bibitem{amaro2017laser}
P.~Amaro-Seoane et~al.,
\emph{Laser Interferometer Space Antenna},
arXiv:1702.00786 (2017),
\doi{10.48550/arXiv.1702.00786}.

\bibitem{ambjorn2005spectral}
J.~Ambj{\o}rn, J.~Jurkiewicz, and R.~Loll,
\emph{The spectral dimension of the universe is scale dependent},
Phys. Rev. Lett., \textbf{95} (2005), 171301,
\doi{10.1103/PhysRevLett.95.171301}.

\bibitem{badurina2020aion}
L.~Badurina et~al.,
\emph{AION: An Atom Interferometer Observatory and Network},
J. Cosmol. Astropart. Phys., \textbf{2020} (2020), 011,
\doi{10.1088/1475-7516/2020/05/011}.

\bibitem{ebadi2021quantum}
S.~Ebadi et~al.,
\emph{Quantum phases of matter on a 256-atom programmable quantum simulator},
Nature, \textbf{595} (2021), 227--232,
\doi{10.1038/s41586-021-03582-4}.

\bibitem{goroff1986ultraviolet}
M.~H. Goroff and A.~Sagnotti,
\emph{The ultraviolet behavior of Einstein gravity},
Nuclear Physics B, \textbf{266} (1986), no.~3-4, 709--736,
\doi{10.1016/0550-3213(86)90193-8}.

\bibitem{hazumi2019litebird}
M.~Hazumi et~al.,
\emph{LiteBIRD: A satellite for the studies of B-mode polarization and Inflation from cosmic background radiation detection},
J. Low Temp. Phys., \textbf{194} (2019), 443--452,
\doi{10.1007/s10909-019-02150-5}.

\bibitem{maldacena2016bound}
J.~Maldacena, S.~H. Shenker, and D.~Stanford,
\emph{A bound on chaos},
J. High Energ. Phys., \textbf{2016} (2016), 106,
\doi{10.1007/JHEP08(2016)106}.

\bibitem{punturo2010einstein}
M.~Punturo et~al.,
\emph{The Einstein Telescope: a third-generation gravitational wave observatory},
Class. Quantum Grav., \textbf{27} (2010), 194002,
\doi{10.1088/0264-9381/27/19/194002}.

\bibitem{reitze2019cosmic}
D.~Reitze et~al.,
\emph{Cosmic Explorer: The U.S. contribution to gravitational-wave astronomy beyond LIGO},
Bull. Am. Astron. Soc., \textbf{51} (2019), 035,
\doi{10.48550/arXiv.1907.04833}.

\bibitem{jacob2008lorentz}
U.~Jacob and T.~Piran,
\emph{Lorentz-violation-induced arrival delays of cosmological particles},
J. Cosmol. Astropart. Phys., \textbf{2008} (2008), 031,
\doi{10.1088/1475-7516/2008/01/031}.

\bibitem{mirshekari2012constraining}
S.~Mirshekari, N.~Yunes, and C.~M. Will,
\emph{Constraining Lorentz-violating, modified dispersion relations with gravitational waves},
Phys. Rev. D, \textbf{85} (2012), 024041,
\doi{10.1103/PhysRevD.85.024041}.

\bibitem{bicepkeck2021constraints}
P.~A.~R. Ade et~al. (BICEP/Keck Collaboration),
\emph{Improved constraints on primordial gravitational waves using Planck, WMAP, and BICEP/Keck observations through the 2018 observing season},
Phys. Rev. Lett., \textbf{127} (2021), 151301,
\doi{10.1103/PhysRevLett.127.151301}.

\bibitem{planck2018parameters}
N.~Aghanim et~al. (Planck Collaboration),
\emph{Planck 2018 results. VI. Cosmological parameters},
Astron. Astrophys., \textbf{641} (2020), A6,
\doi{10.1051/0004-6361/201833910}.

\bibitem{silvafilho2026flows} R.~M. Silva-Filho, \emph{Geometric Flows, Partial Differential Equations, and Variational Dynamics on Matrix and Tensor Manifolds}, Chapter~2 of \emph{Geometry, Tensors, and Quantum Gravity}, Universidade Federal de Lavras, 2026.

\bibitem{silvafilho2026functional} R.~M. Silva-Filho, \emph{Functional Realizations of Matrices and Hypertensors: Emergent Invariants and Geometric Measures}, Chapter~1 of \emph{Geometry, Tensors, and Quantum Gravity}, Universidade Federal de Lavras, 2026.

\bibitem{silvafilho2026grand} R.~M. Silva-Filho, \emph{Emergent Spacetime, Holographic Entanglement, and Quantum Geometry: Proved Results, Conjectures, and Observational Limits}, Chapter~12 of \emph{Geometry, Tensors, and Quantum Gravity}, Universidade Federal de Lavras, 2026.

\bibitem{silvafilho2026noneuclidean} R.~M. Silva-Filho, \emph{Minimax Extrinsic Curvature in Non-Euclidean Geometries and General Relativity: Variational Theory, Space Forms, and ADM Spacetime Slicings}, Chapter~8 of \emph{Geometry, Tensors, and Quantum Gravity}, Universidade Federal de Lavras, 2026.

\bibitem{silvafilho2026pascal} R.~M. Silva-Filho, \emph{Analytic Continuation of Pascal's Simplex: Continuous Multinomial Integrals, Polytope Boundary Recurrences, and Simplicial Fractional Calculus}, Chapter~3 of \emph{Geometry, Tensors, and Quantum Gravity}, Universidade Federal de Lavras, 2026.

\bibitem{silvafilho2026spacetime} R.~M. Silva-Filho, \emph{Emergent Spacetime and Quantum Geometry: Tensor Networks, Holonomies, and Entanglement (Established Results and Conjectures)}, Chapter~11 of \emph{Geometry, Tensors, and Quantum Gravity}, Universidade Federal de Lavras, 2026.

\end{thebibliography}

\end{document}
